\section{Contracts, verifiers and audit}
\label{sec:method}
\paragraph{Requirement contracts.} A contract is a set of requirements, each with a statement, a kind (change, invariant, constraint, non-goal), a status (explicit, inferred, superseded, unresolved), the event that introduced it, and evidence spans that must match the session text exactly. Proposed requirements whose quotes cannot be aligned to user messages are rejected (or downgraded to \emph{inferred} if a rationale is given). Supersession links must point forward in time; active requirements may not conflict; unresolved requirements cannot be scored; inferred requirements need a human review before a task is accepted. The \emph{initial-only} contract is recovered by the same proposer from the first user request alone; the \emph{full} contract from the whole session.

\paragraph{Verifier conditions.} From a contract, a generator writes 6--14 behavioural tests (C1 initial/plain, C3 full/plain). The audited conditions start from these tests and apply the audit below (C2 initial/audited, C4 full/audited). Baselines: B1, the repository's tests at the reference state restricted to those passing on the reference; B2, a generator given the whole session transcript without contract machinery. All generators see the base repository and the reference diff and use the same model and prompt template.

\paragraph{Audit (audit:v1).} Using development candidates only (the reference solution repeated three times, the no-op, four development attacks that each violate one contract requirement, and one development valid alternative), the audit requires: the reference passes on every repeat, the no-op and every attack fail, the alternative passes, every scored requirement is mapped to a check, and static lint is clean. At most three revision rounds are allowed; the verifier is then frozen.

\paragraph{Process-separated grading.} Hidden tests run in a trusted grader process; candidate code runs only in child processes under an unprivileged user that cannot read the tests or write the results. Each child's result is authenticated with a per-call key; the host recomputes the outcome from the results file and a per-run nonce, so missing or stale results are verifier errors, never passes. Containers run without network, with all capabilities dropped except those needed to switch users, and with resource limits. Candidate edits to test files and test configuration are discarded before grading. \S\ref{sec:repro} lists the isolation probes.
